\section{Security}\label{sec:offline-security}

%%%% proof 1 - the generic construction securely realises the ideal
%functionality

\subsection{Proof
of~\Cref{thm:offline-construction-security}}\label{sec:offline-proof-1}

Let $\algo{prot}$ denote the protocol of Protocol~\ref{prot:offline}. We write
$\algo{prot}$ rather than $\Pi$ for the protocol, as $\Pi$ already denotes the
withdrawal proof. Let $\adv$ be a real world adversary and let $\simu$ be an
ideal world simulator that interacts with system participants through ideal
functionality $\idfu{\algo{op}}$. We show that
$\textsc{Exec}_{\algo{prot},\adv}(\secparam)$ is computationally
indistinguishable from $\textsc{Exec}_{\idfu{\algo{op}},\simu}(\secparam)$, in
the standalone sense of~\Cref{sec:ideal-func}. We define a sequence of
indistinguishable games that start with Game 0 representing the real world, and
end with Game 12 representing the ideal world. Throughout, we call a user
corrupt if both their device and secure element are compromised. The more
general split corruption setting follows analogously. We also simplify
exposition by treating the randomness $\rho$ produced by $\idfu{\algo{op}}$
during $\OPwithdraw$ and $\OPexecute$ as the output of $\com(\ck, \cdot;
\cdot)$.

We rely on the following security properties of the underlying primitives:

\begin{enumerate}

  \item[(P1)] The commitment scheme is binding: no PPT adversary can find two
openings of the same commitment to distinct messages.

  \item[(P2)] The encryption scheme is IND-CPA secure: ciphertexts are
computationally indistinguishable from random ones for any PPT adversary.

  \item[(P3)] The blind signature scheme used by $\algo{cb}$ is EUF-CMA secure:
no PPT adversary can produce a valid signature without querying the signer.

  \item[(P4)] The credential signature scheme used by $\reg$ is EUF-CMA secure.

  \item[(P5)] The zero knowledge proofs and signatures of knowledge, including
the joint proof of knowledge, are simulation extractable in the random oracle
model: a valid witness can be extracted from any accepted proof, even by an
extractor that has itself produced simulated proofs, and simulated proofs are
indistinguishable from real ones.

  \item[(P6)] The withdrawal proof $\Pi$ is online extractable: its witness can
be recovered without rewinding the prover. This is needed only where $\simu$
extracts in order to continue the simulation rather than to derive a
contradiction, which is the case for $\Pi$ alone.

\end{enumerate}

\paragraph{The simulator.} Before giving the games we describe $\simu$ in one
place. $\simu$ runs $\algo{prot}$ internally on behalf of every honest party,
and relays between $\adv$ and $\idfu{\algo{op}}$. It maintains two maps,
$\mathbb{R}$ from sequences of pseudonyms to the token identifiers
$\idfu{\algo{op}}$ uses, and $\mathbb{C}$ for the inverse direction, so that it
can translate between the transcripts $\adv$ sees and the state
$\idfu{\algo{op}}$ holds. For honest parties $\simu$ never has a witness to
prove, so every proof it produces is simulated and every commitment and audit
ciphertext it produces is random; this is what makes its transcripts independent
of the identities and token identifiers it does not learn. For corrupt parties
$\simu$ has no witness either, so it extracts one from the proofs $\adv$
produces: from $\Pi$ at withdrawal, in order to learn the randomness that fixes
the token identifier and so continue the simulation, and from $\Sigma$,
$\sigma_i$, and $\Gamma_i$ elsewhere, in order to exhibit a forgery or a double
opening and terminate. The first of these is why (P6) is needed and the second
is why rewinding suffices for the rest.

\begin{proof}

\paragraph{Game 0.} The real world execution of $\algo{prot}$.

\paragraph{Game 1.} As Game 0, except $\simu$ runs $\algo{prot}$ internally on
behalf of all honest parties, forwarding inputs and distributing outputs. This
change is purely syntactic, so Game 1 is identically distributed to Game 0.

\paragraph{Game 2.} As Game 1, except $\simu$ introduces a dummy
$\idfu{\algo{op}}$ as a communication relay between honest parties, who now
communicate via $\idfu{\algo{op}}$ rather than directly. This change is again
syntactic, so Game 2 is identically distributed to Game 1.

\paragraph{Game 3.} As Game 2, except $\idfu{\algo{op}}$ handles $\OPcreate$
queries from honest users. When an honest user $\algo{usr}$ initialises an
account of value $v$ with intermediary $\algo{Int}$, $\simu$ calls
$\idfu{\algo{op}}$ with $(\OPcreate, \algo{usr}, v)$ in place of running
$\algo{prot}$. $\algo{prot}$ always accepts when the user has no prior account
and $\idfu{\algo{op}}$ performs the same check, so the two are identical.

\paragraph{Game 4.} As Game 3, except corrupt users may also initiate account
creation. $\simu$ calls $\idfu{\algo{op}}$ and mirrors any abort to
$\algo{Int}$. $\idfu{\algo{op}}$ aborts exactly when $\algo{prot}$ does, on
attempts to register duplicate accounts, so no new distinguishing event arises.

\paragraph{Game 5.} As Game 4, except $\idfu{\algo{op}}$ handles $\OPregister$
queries: $\simu$ calls $(\OPregister, \algo{usr}, \algo{se}, \algo{dev})$
whenever $\reg$ receives a valid registration request.  $\idfu{\algo{op}}$
rejects only duplicate registrations, matching $\algo{prot}$, so the games are
indistinguishable.

\paragraph{Game 6.} As Game 5, except $\idfu{\algo{op}}$ handles $\OPwithdraw$
queries from honest users. When an honest $\algo{usr}$ withdraws value $v$,
$\simu$ calls $(\OPwithdraw, \algo{Int}, v)$ and learns $(\algo{usr}, v)$. It
then produces a simulated withdrawal transcript. In it, the request commitment
$C$ is a uniformly random commitment, the two blind signing ciphertexts
$(\hat{\ciph}_0, \hat{\ciph}_2)$ are random ciphertexts under a freshly
generated key, and $\Pi$ is a simulated proof. Note that these ciphertexts are
under the withdrawer's own key and carry the hidden messages to the signer; they
are not the audit ciphertexts of a payment, which are under $\pk_\auditor$ and
appear first in Game 8. By (P1), (P2), and (P5), this transcript is
computationally indistinguishable from one produced by $\algo{prot}$.

\paragraph{Game 7.} As Game 6, except $\idfu{\algo{op}}$ handles $\OPwithdraw$
queries from corrupt users. When $\algo{usr}$ submits a withdrawal, $\simu$
extracts the randomness $s$ from the attached proof using (P6), computes $\rho
\leftarrow \com(\ck, (\algo{uid}, v); s)$, and calls $(\OPwithdraw, \algo{Int},
v)$ with $\rho$. Both $\idfu{\algo{op}}$ and $\algo{prot}$ reject if and only if
the request commitment $C$ to $(\algo{uid}, v, s)$ is a duplicate, which would
imply a duplicate $\rho$, so the accounting checks agree and the games are
indistinguishable.

\paragraph{Game 8.} As Game 7, except $\idfu{\algo{op}}$ handles $\OPexecute$
queries from honest payers. We first deal with the first payment after
withdrawal, the case $i = 0$, then generalise.

Case $i = 0$. $\algo{usr}$ pays $\algo{usr}'$ by calling $(\OPexecute, \rho, v,
\algo{usr}')$.

\begin{itemize}

  \item If $\algo{usr}'$ is honest: $\simu$ learns $(\OPexecute, v, \rho,
\rho')$ from $\idfu{\algo{op}}$, and produces a simulated payment $p_0 = (v,
\com_0, \com_1, \ciph_0, \sigma_0, \Gamma_0, \Sigma)$ with random commitments
$\com_b \leftarrow \com(\ck, r_b; s_b)$ for fresh $r_b, s_b$, a random
ciphertext $\ciph_0$, and simulated proofs $(\sigma_0, \Gamma_0, \Sigma)$,
possible by simulation extractable security of the NIZK system (P5).

  \item If $\algo{usr}'$ is corrupt: $\simu$ additionally learns the identities
of both parties. It computes $\rho' \leftarrow \com(\ck, \algo{uid}'; s')$ from
the commitment $\com_1$ that $\algo{usr}'$ supplies in $\algo{prot}$, calls
$\idfu{\algo{op}}$, and sets $\ciph_0 \leftarrow \PKEencrypt(\pk_\auditor,
\algo{uid}; \rho_0)$.

\end{itemize}

In both cases $\simu$ records $\mathbb{R}[(\com_0, \com_1)] \gets (\rho, \rho')$
and $\mathbb{C}[(\rho, \rho')] \gets (\com_0, \com_1)$.

Case $i > 0$. $\simu$ recovers $(\com_0, \ldots, \com_i)$ from
$\mathbb{C}[\vec{\rho}]$, produces a simulated payment $p_i$ with simulated
proofs and random ciphertexts (or $\ciph_i \leftarrow \PKEencrypt(\pk_\auditor,
\algo{uid}_i; \rho_i)$ if $\algo{usr}'$ is corrupt), and updates the maps as
needed. By IND-CPA security of the encryption scheme (P2) and simulation
extractable security of the NIZK scheme (P5), simulated payments are
indistinguishable from real ones, so Game 8 is indistinguishable from Game 7.

\paragraph{Game 9.} As Game 8, except $\idfu{\algo{op}}$ also processes
$\OPexecute$ queries from corrupt payers. We show that if $\idfu{\algo{op}}$
rejects a payment, an honest $\algo{usr}'$ in $\algo{prot}$ also rejects it,
again showing the case $i = 0$ before generalising.

Let $p_0 = (v, \com_0, \com_1, \ciph_0, \sigma_0, \Gamma_0, \algo{hist}_0)$ be
the payment produced by corrupt $\algo{usr}$. By (P5), $\simu$ extracts
$(\algo{uid}, s)$ from $\Sigma$, computes $\rho \leftarrow \com(\ck,
(\algo{uid}, v); s)$, and calls $(\OPexecute, \rho, v, \algo{usr}')$.

$\algo{usr}'$ honest. If $\algo{usr}'$ accepts $p_0$ but $\idfu{\algo{op}}$
rejects, exactly one of the following holds:

\begin{enumerate}[label=\alph*)]

  \item $\algo{tokens}[\rho] = \bot$: no prior $\OPwithdraw$ matches $\rho$. By
(P5), $\simu$ extracts a blind signature $\psi$ and message triple $(\algo{uid},
v, s)$ from $\Sigma$. Since $\algo{cb}$ never issued $\psi$, no corresponding
withdrawal exists, so outputting $(\psi, \algo{uid}, v, s)$ breaks EUF-CMA of
the blind signature scheme (P3).

  \item $\algo{tokens}[\rho] = \langle \algo{usr}, v^* \rangle$ with $v^* \neq
v$: we can extract $(\ck, \algo{uid}', s', \com_0)$ from $\Sigma$ such that
$\ComOpen(\ck, \algo{uid}', s', \com_0) = 1$, giving $\rho \leftarrow \com(\ck,
\algo{uid}'; s')$ with value $v$. Combined with $\rho \leftarrow \com(\ck,
\algo{uid}^*; s^*)$ with value $v^*$ from the prior withdrawal, $v \neq v^*$
yields two distinct openings of $\rho$, contradicting binding of the commitment
scheme (P1).

  \item $\algo{tokens}[\rho] = \langle \algo{usr}^*, v \rangle$ with
$\algo{usr}^* \neq \algo{usr}$: if $\algo{usr}^*$ is corrupt, the same two
representations of $\rho$ as in (b) give $\algo{uid} \neq \algo{uid}^*$, again
contradicting (P1). If $\algo{usr}^*$ is honest, $\com_0 \leftarrow \com(\ck,
\algo{uid}^*; s^*)$ was chosen by $\simu$, so $\algo{usr}$ cannot have guessed
it except with negligible probability. If $\algo{usr}'$ accepts $p_0$, then
$\sigma_0$ is a valid signature of knowledge proving possession of a credential
on $(\algo{uid}, \sk)$, and by (P5) $\simu$ extracts credential $\varphi$ and
pair $(\algo{uid}, \sk)$. Since $\algo{uid}$ was assigned randomly by $\simu$
and was not registered, this gives a forgery of the credential scheme (P4).

  \item $\algo{usr}$ is not registered: if $\algo{usr}'$ accepts, $\sigma_0$
verifies, and by (P5) $\simu$ extracts a credential forgery $(\varphi,
(\algo{uid}, \sk))$, contradicting (P4).

\end{enumerate}

$\algo{usr}'$ corrupt. $\simu$ sends $\algo{accept}$ to $\idfu{\algo{op}}$ on
$\algo{usr}'$'s behalf, provides $\rho' \leftarrow \com(\ck, (\algo{uid}', v);
s')$ from $\com_1$, and extends the maps.

$i > 0$. Suppose $\mathbb{R}[(\com_0, \ldots, \com_i)]$ is not yet recorded.
$\simu$ fills in the missing intermediate payments as follows. It extracts
participant identifiers from the proof $\Gamma_j$ at the first unrecorded index
$j$ using (P5), calls $\OPexecute$ for each, and extends the maps until $j = i -
1$. These extractions are on instances already fixed by the transcript $p_i$
that $\algo{usr}$ has sent. A corrupt user cannot make a later instance depend
on the extraction of an earlier one, so the extractions are independent rather
than nested: $\simu$ performs them sequentially, each by rewinding to the same
transcript, and the forking loss is linear rather than exponential in the length
of the chain, which the epoch cap of~\Cref{sec:offline-construction} bounds in
any case. The cases above still apply, replacing $p_0$ with $p_i$ and $\rho$
with $(\rho_0, \ldots, \rho_i)$.  Since the probability that an honest
$\algo{usr}'$ accepts a payment that $\idfu{\algo{op}}$ rejects is negligible
under (P1)-(P6), this game is indistinguishable from Game 8.

\paragraph{Game 10.} As Game 9, except $\idfu{\algo{op}}$ handles $\OPdeposit$
queries from honest users. $\simu$ recovers $(\com_0, \ldots, \com_n)$ from
$\mathbb{C}[\vec{\rho}]$ and produces a simulated deposit $(\tau_n, \ciph_n,
\sigma_n, \Gamma_n, \algo{hist}_n)$ with simulated proofs and random
ciphertexts. $\idfu{\algo{op}}$ rejects a deposit only when the token does not
belong to $\algo{usr}$ or has been previously deposited. An honest user never
deposits a token they do not own, so the only case requiring argument is double
spend detection.

Each $\rho_i$ is determined by $\com_i$ via $\rho_i \leftarrow \com(\ck,
(\algo{uid}_i, v); s_i)$, so the mapping $\rho_i \leftrightarrow \com_i$ is
one-to-one by binding of the commitment scheme (P1). $\rho_0$ is fixed by
$\com_0$, which is fixed by the blind signature issued by $\algo{cb}$, which
binds $(\algo{uid}_0, v, s_0)$. Two withdrawals yielding the same $\com_0$ would
have to agree on $\algo{uid}_0$ and $s_0$, and $s_0$ is sampled uniformly by the
secure element, so this occurs with negligible probability. A corrupt user
cannot alter $\com_0$, so any two deposits sharing $\vec{\rho}[0]$ are double
spends and will be detected by $\idfu{\algo{op}}$, so Game 10 is
indistinguishable from Game 9.

\paragraph{Game 11.} As Game 10, except $\idfu{\algo{op}}$ also handles
$\OPdeposit$ queries from corrupt users. When $\algo{tokens}[\vec{\rho}] \neq
\langle \algo{usr}, v \rangle$, exactly one of the following holds. Each is
handled by the same argument as in Game 9, with the payment in question replaced
by the deposit in question.

\begin{enumerate}[label=\alph*)]

  \item $\algo{tokens}[\vec{\rho}] = \bot$: extracting $\Sigma$ via simulation
extractability of the NIZK scheme (P5) gives a blind signature forgery,
contradicting unforgeability of the signature scheme.

  \item $\algo{tokens}[\vec{\rho}] = \langle \algo{usr}^*, v^* \rangle$ with
$v^* \neq v$: two valid openings of $\rho_0$ under $\ComOpen$ with distinct
values violate binding of the commitment scheme (P1).

  \item $\algo{tokens}[\vec{\rho}] = \langle \algo{usr}^*, v \rangle$ with
$\algo{usr}^* \neq \algo{usr}$: either binding of the commitment scheme (P1) is
violated, or extracting $\sigma_0$ via (P5) yields a credential forgery,
contradicting EUF-CMA of the credential signature scheme (P4).

  \item $\algo{usr}$ unregistered: $\sigma_0$ verifying implies a credential
forgery, contradicting unforgeability of the credential signature scheme.

\end{enumerate}

If the token has been transferred multiple times before deposit, $\simu$ fills
in unrecorded payments using $\Gamma_j$ as in Game 9, and the same cases apply
inductively. If another deposit chain sharing $\vec{\rho}[0]$ has already been
recorded, the same argument as in Game 10 shows the initial commitment cannot be
altered, so Game 11 is indistinguishable from Game 10.

\paragraph{Game 12.} As Game 11, except $\idfu{\algo{op}}$ handles $\OPidentify$
queries. The auditor calls $(\OPidentify, \rho)$ and expects the set of double
spenders $\algo{DS}$. We show that the outputs of $\idfu{\algo{op}}$ and
$\algo{prot}$ agree except with negligible probability.

In $\algo{prot}$, the double spender at index $j$ is identified by decrypting
$\ciph_j$, where $j$ is the last index at which the two deposit chains agree.
Game 11 fixes the commitment $\com_j$, by binding of the commitment scheme and
EUF-CMA of the blind signature scheme. This carries through the chain via
simulation extractability of $\sigma_{j-1}$, which is taken over the message
$(v, \com_{j-1}, \com_j)$ and so fixes $\com_j$ as the destination of the
payment that delivered the token. The ciphertext $\ciph_j \leftarrow
\PKEencrypt(\pk_\auditor, \algo{uid}_j; \rho_j)$ encrypts the same value as is
committed in $\com_j$, as certified by proof $\Gamma_j$. Were this false, it
would either yield two valid openings of $\com_j$ under $\ComOpen$,
contradicting binding (P1), or a proof $\Gamma_j$ with no valid witness,
contradicting simulation extractability (P5).

An honest user cannot be framed. To have an honest $\algo{uid}^*$ returned, the
adversary would need $\ciph_j$ to decrypt to $\algo{uid}^*$ while $\com_j$ is
the pseudonym recorded at index $j$. If $\com_j$ was generated by an honest
payee, then producing $\Gamma_j$ requires opening it, and by (P5) $\simu$
extracts that opening, contradicting (P1) since the honest payee's own opening
is different, or contradicting hiding of $\com_j$ since $\algo{uid}^*$ was drawn
at random. If $\com_j$ was instead supplied by the adversary as a commitment to
$\algo{uid}^*$, then spending the token at index $j$ also requires $\sigma_j$, a
signature of knowledge of a credential on $\algo{uid}^*$, from which (P5)
extracts a forgery contradicting (P4). The same argument covers $j = 0$, where
$\com_0$ is chosen by the withdrawer rather than by a payee.

Suppose $\ciph_j$ decrypts to an unregistered user. Then $\sigma_j$ is a valid
signature of knowledge over a credential for that identifier. By simulation
extractability of the NIZK (P5), this contradicts EUF-CMA of the credential
scheme (P4). All distinguishing events occur with negligible probability, so
Game 12, which is the ideal world, is indistinguishable from Game 11. This
completes the proof.

\end{proof}

%%%% proof 2 - the concrete instantiation, corollary of Theorem 1

\subsection{Proof
of~\Cref{thm:offline-instantiation-security}}\label{sec:offline-proof-2}

The instantiation of~\Cref{sec:offline-full-inst} is an instance of the
construction that sets the commitment scheme to Pedersen commitments in
$\group{1}$, the credential signature scheme to BBS+, the blind signature scheme
to Pointcheval-Sanders signatures over committed inputs, the public key
encryption scheme to ElGamal in $\group{1}$, and the joint proof of knowledge to
Primitive~\ref{prim:bbs-joint-proof}. It therefore suffices to discharge (P1) to
(P6).

\begin{proof}

\begin{description}

  \item[(P1) The commitment scheme is binding.] Pedersen commitments are
computationally binding under the discrete logarithm assumption in $\group{1}$.
The generators are derived by hashing into $\group{1}$, as described
in~\Cref{sec:offline-instantiation}, so no participant, including the central
bank, knows the discrete logarithm relations that would break binding.

  \item[(P2) The encryption scheme is IND-CPA secure.] ElGamal in $\group{1}$ is
IND-CPA secure under the DDH assumption in $\group{1}$.

  \item[(P3) The blind signature scheme is EUF-CMA secure.] Pointcheval-Sanders
signatures are existentially unforgeable under the $q$-MSDH assumption of
Assumption~\ref{assum:qmsdh}~\cite{RSA:PoiSan18}, and blind signing over
committed inputs preserves unforgeability~\cite{NDSS:SABMD19}.

  \item[(P4) The credential signature scheme is EUF-CMA secure.] BBS+ signatures
are existentially unforgeable under the $q$-SDH
assumption~\cite{SCN:AuSusMu06,EPRINT:CamDriLeh16}.

  \item[(P5) The proofs and signatures of knowledge are simulation extractable.]
Every proof in the instantiation is a sigma protocol over discrete logarithms
with the Fiat-Shamir heuristic applied~\cite{C:FiaSha86}. Such proofs have
quasi-unique responses over a prime order group, and are therefore simulation
extractable in the random oracle model~\cite{INDOCRYPT:FKMV12}. This includes
the joint proof of Primitive~\ref{prim:bbs-joint-proof}, which is a conjunction
of two such protocols under separate challenges, so the same argument applies to
each conjunct and hence to the whole. We rely on this rather than on the
soundness established for the primitive in~\cite{USENIX:HesSinSor23}, which is
weaker: $\simu$ must extract from proofs produced by $\adv$ while itself
producing simulated proofs for honest users, and plain soundness does not cover
that. The resulting extractor yields the payer's identifier and credential from
an accepted $\sigma_i$.

  \item[(P6) The withdrawal proof is online extractable.] The extractor
of~\cite{INDOCRYPT:FKMV12} proceeds by rewinding, which suffices wherever
$\simu$ extracts once in order to derive a contradiction. In Game 7, however,
$\simu$ must extract from $\Pi$ in order to continue the simulation. It does so
on commitments a corrupt user chooses adaptively and without bound, and the cost
of rewinding compounds across such extractions.  We therefore require $\Pi$ to
admit an online extractor, as in~\cite{USENIX:HesSinSor23} for the analogous
proof accompanying blind issuance. This applies to $\Pi$ alone. It is not
required of $\sigma_i$ or $\Gamma_i$, so it bears on the mobile device during
withdrawal rather than on the secure element or on any part of a payment.

\end{description}

Because $\simu$ rewinds to extract, the argument is a standalone one, in the
sense fixed in~\Cref{sec:ideal-func}, and we do not claim universal
composability. Discharging the remaining rewinding extractions with an online
extractor, as (P6) does for $\Pi$, would give a straight line simulator and with
it a composable result, at the cost of a larger proof at each hop. We leave that
trade to future work.

The instantiated protocol therefore securely realises ideal functionality
$\idfu{\algo{op}}$ in the random oracle model.

\end{proof}
